\section{五八九年之后：接收不是把差异消除}
\label{sec:synthesis-reunification-afterlives}

陈亡的消息沿江水与州郡传开时，隋朝已经赢得了最高政治层面的决定性胜利。建康受降、南朝皇统结束和隋的天下名义，可以明确写在五八九年的年序中；这并不意味着江南的户籍、寺院、家族、市场和水路在同一天变成北方的复制品。再统一值得被写成转折，正因为它把长期敌对的朝廷关系改造成接收、登记、运输和裁判的任务，而非因为它使这些任务立刻完成。

\subsection{受降的日期与地方的时间}

受降仪式能够宣布谁拥有合法命令权，不能逐项清点一地的仓储、一县的簿册或一户的迁居。新朝接收州郡时，必须决定哪些官员留用，哪些旧有税役继续，哪些军队与船只重新编入，哪些寺院资源需要承认或清查。地方官与居民也要判断新命令是否会持续，旧有债务、婚姻与依附关系能否被承认。政治中心的更替因而只是地方时间的开始，而不是终点。

这一差别同样适用于城市。建康并未因为陈亡而失去长江下游的道路、市场、寺院和熟练劳力；新政权若要利用它，就必须接收而非简单清空这些关系。江陵、荆州及江淮各地的水路网络也仍受季节、船户、堤防和粮食来源约束。隋能在较大范围内协调资源，是统一的重要收益；协调是否能稳定，则取决于能否让先前服务于南北竞争的基础设施转为共同使用。

\subsection{北方遗产与南方遗产并不对称}

隋继承北周的关中军政框架、对河北资源的接收经验以及同草原势力周旋的条件，这些都是南下得以实现的前提。它并非从一张空白地图上设计统一。与此同时，陈及更早南朝留下的州郡、文书、寺院、家族与经济网络，也不是等待被动“整合”的遗物。南北两种政治经验在新的政府中相遇，既提供可用的人员和制度，也带来不同的地方习惯与利益。

因此，不宜把五八九年写成“北方制度压倒南方社会”，或把南方写成完全不受北方改变的连续世界。统一后的行政安排会重新塑造户籍、赋役和官员流动，南方的水路、城市和家族关系也会限制这种塑造能够多快、以何种方式完成。历史解释应当让双向作用保持可见：国家能以任命、征收与道路改变地方，地方同样以资源、知识和既有关系改变国家的执行。

\subsection{宗教与书写的接收}

寺院和经卷显示接收并不限于官府。北周末年的宗教政策、隋初的调整和南方既有寺院网络，使新政权需要面对人员、建筑、供养和文本的复杂遗产。赞助或清查都不是纯粹象征：它们会影响僧侣能否居住、物资怎样处置、城市中哪些组织得到保护。但没有材料允许我们断言每座寺、每位僧尼或每位信众同时作出相同反应。再统一之后的宗教史必须保留政策宣布与地方实践的距离。

书写同样跨过朝代。官府需要重新组织档案与簿册，家族继续保存墓志、家训和藏书，寺院与读写者使经卷流通。隋的政权名义能够为这些活动提供新的框架，却不能倒转先前的抄写、迁徙和丢失。对读者而言，这提醒我们：所谓“连续性”并非一切不变，而是人们不断在新条件下重建可保存、可阅读、可传递的东西。

\begin{turningpoint}[再统一的制度得失：把接收变成可持续的治理]
隋的统一减少了南北对峙造成的最高政治边界，使交通、征发和任命有机会在更大范围协调；成本是新政府必须承担更多区域差异，并把接收的风险带入自己的户籍、财政和边防体系。制度是否成功，不只看受降是否完成，还看地方能否获得较可预期的登记、供给和裁判。
\end{turningpoint}

\subsection{本卷如何承接下一段历史}

两晋南北朝没有留下一份可以直接交给隋的完整答案。它留下的是经由多次统一与崩解检验过的问题：流动人口如何安置，军粮如何跨越河山，登记如何经过地方中介，寺院和家族怎样拥有资源，边地怎样连接都城，文本怎样在失都之后保存记忆。隋能够利用其中的经验，也必须承担其中未解的压力。下一阶段的制度规模、交通计划与帝国野心，正是在这些旧问题仍未消失的条件下展开。

读者若从这里回望二八〇年的西晋统一，会发现两端之间并不是“统一—分裂—统一”的三格图。每一次统一都把更广的道路、人群与资源纳入同一责任；每一次崩解都让这些责任重新落到地方、家庭和中介身上。将五八九年写成开放的结尾，既承认隋的胜利，也保留那些继续穿过王朝更替的生活与材料。
